% 三3正承担.tex - 正在承担的项目

申请人目前承担以下与本项目相关的国家级研究项目：

\begin{table}[htbp]
\centering
\small
\begin{tabular}{|p{2.5cm}|p{4.8cm}|p{2.3cm}|p{1.8cm}|p{2.8cm}|}
\hline
\textbf{资助机构} & \textbf{项目名称} & \textbf{批准号} & \textbf{起止时间} & \textbf{与本项目关系} \\
\hline
国家自然科学基金委（青年基金） & 基于里德堡原子阵列的可并行量子算法研究 & 12404568 & 2025.01--2027.12 & 提供里德堡算法平台理论基础 \\
\hline
科技部（国家重点研发计划，骨干成员） & 量子信息处理的数学基础——非交换数学理论的发展与应用 & 2024YFE0102500 & 2024.06--2027.05 & 提供形式化方法理论支撑 \\
\hline
\end{tabular}
\end{table}

\subsection*{在研项目与本项目的关系}

上述在研项目从不同角度与本项目相关：

\textbf{青年基金项目（12404568）}聚焦于里德堡原子阵列的可并行算法设计，已完成三角晶格编码方案的理论设计与数值验证。该项目为本项目提供了问题映射的核心方法论基础，本项目将在此基础上发展机器辅助的形式化证明框架，实现从经验性设计到可验证正确性的跨越。
\textbf{国家重点研发计划项目（2024YFE0102500）}研究量子信息处理的数学基础，特别是非交换代数结构的形式化表示。该项目的形式化方法研究经验为本项目基于整数规划的正确性验证提供理论参考，特别是在处理组合优化问题的约束结构与对称性方面。

本项目与上述在研项目在研究对象上有关联，但科学目标和技术路线独立。
